(1) 取\(\beta=1/(k_BT)\)，为防止场项重复计数，将每条键的Boltzmann因子定义为\(T_{s,s'}=\exp[\beta Jss'+(\beta h/2)(s+s')]\)。用基底\((+1,-1)\)得到\(\mathsf T=\begin{pmatrix}e^{\beta(J+h)}&e^{-\beta J}\\e^{-\beta J}&e^{\beta(J-h)}\end{pmatrix}\)；\(\operatorname{tr}\mathsf T=2e^{\beta J}\cosh\beta h\)，行列式\(e^{2\beta J}-e^{-2\beta J}\)，故\(\lambda_\pm=e^{\beta J}[\cosh(\beta h)\pm\sqrt{\sinh^2(\beta h)+e^{-4\beta J}}]\)。
(2) 周期链配分函数\(Z_N=\operatorname{tr}\mathsf T^N=\lambda_+^N+\lambda_-^N\)。因\(T>0\)和\(J>0\)，最大本征值\(\lambda_+\)严格占优，故\(f=\lim_{N\to\infty}(-k_BT/N)\ln Z_N=-k_BT\ln\lambda_+\)。每自旋磁化\(m=-\partial f/\partial h=\sinh(\beta h)/\sqrt{\sinh^2(\beta h)+e^{-4\beta J}}\)；它的符号与外场同向且满足\(|m|\le1\)。
(3) 在\(h=0\)，\(\lambda_+=2\cosh(\beta J)\)、\(\lambda_-=2\sinh(\beta J)\)。插入两个\(s\)算符或用最近邻Markov条件概率，得到\(\langle s_0s_r\rangle=(\lambda_-/\lambda_+)^r=[\tanh(\beta J)]^r=e^{-r/\xi}\)，\(\xi=-1/\ln\tanh(\beta J)\)。从磁化表达式求\(\chi=\partial m/\partial h|_{h=0}=\beta e^{2\beta J}\)；它也等于\(\beta[1+2\sum_{r\ge1}(\tanh\beta J)^r]=\beta(1+\tanh\beta J)/(1-\tanh\beta J)\)，利用双曲恒等式二者完全相等。
(4) 对任何有限\(\beta\)，\(e^{-4\beta J}>0\)，所以\(\lim_{h\to0}m(h)=0\)，不存在\(T>0\)下的自发有序相；\(T\to0\)时\(\xi\sim\tfrac12e^{2\beta J}\)发散、\(\chi\sim\beta e^{2\beta J}\)发散。这一“零温临界”表明\(T\downarrow0\)与\(h\to0\)的极限不可随意交换，也展示了有限程一维Ising模型与高维临界行为的不同。
